\documentclass[a4paper,11pt]{article}


\usepackage{amssymb,amsmath}
\usepackage[frenchb]{babel}
\usepackage[T1]{fontenc}
\usepackage[utf8]{inputenc}
\usepackage{hhline}
\usepackage{enumerate}
\usepackage{a4wide}


\newcommand{\T}{\mathbb{T}}
\newcommand{\N}{\mathbb{N}}
\newcommand{\Z}{\mathbb{Z}}
\newcommand{\Q}{\mathbb{Q}}
\newcommand{\R}{\mathbb{R}}
\newcommand{\C}{\mathbb{C}}
\newcommand{\K}{\mathbb{K}}


\def\om{\omega}
\def\eps{\varepsilon}
\def\DR{\mathcal{C}^\infty_c(\R)}

\DeclareMathOperator{\supp}{supp}

% lignes epaisses
\def\ligne#1{\leaders\hrule height #1\linethickness \hfill}

%-------------------environnements-----------------------------
\newcount\numexo
\numexo=1

\newenvironment{exercice}{\noindent\textbf{Exercice \number\numexo.---}}{\global\advance\numexo by 1}
\newenvironment{pb}{\noindent\textbf{Probl\`eme}}


\frenchspacing



%\usepackage{lmodern}
%
\usepackage{scrextend}
%%%%%%%%%%%%%%%%%%%%%%%%%%%%%%

%
\usepackage{helvet}
\renewcommand{\familydefault}{\sfdefault}
%\usepackage{rotating}
%
\thispagestyle{empty}
%
%\usepackage{inconsolata}
%
%
%\newcommand{\bigsize}{\fontsize{16pt}{20pt}\selectfont}

\usepackage{setspace}
\onehalfspacing


%%%%%%%%%%%%%%%%%%%%%%%%%%%%%%%%%%%%%%%%%%%%%%%%%%%%%%
% quelques symboles utiles en mathematiques
%%%%%%%%%%%%%%%%%%%%%%%%%%%%%%%%%%%%%%%%%%%%%%%%%%%%%%


%\newcommand{\R}{\mathbb{R}}
%\newcommand{\PP}{\mathbb{P}}
%\newcommand{\E}{\mathbb{E}}
%\newcommand{\N}{\mathbb{N}}
%\newcommand{\Q}{\mathbb{Q}}
%\newcommand{\Z}{\mathbb{Z}}

\newcommand{\sh}{{\mathop{\rm sh}}}
\newcommand{\ch}{{\mathop{\rm ch}}}
%\newcommand{\th}{{\mathop{\rm th}}}%commande deja d馭inie
\newcommand{\var}{\mathop{\rm Var}}
\newcommand{\cov}{\mathop{\rm Cov}}

                



\oddsidemargin  -0.25cm
\evensidemargin -0.25cm
\textwidth      17cm
%\headheight     0.in
%\headheight     -.25in
\topmargin      -2.5cm
\textheight     27cm





%%%%%%%%%%%%%%%%%%%%%%%%%%%%%%%%%%%%%%%%%%%%%%%%%


% les lettres grecques
\def\a{\alpha}
\def\b{\beta}
\def\l{\lambda}
%\newcommand{\a}{\alpha}\newcommand{\b}{\beta}\newcommand{\l}{\lambda}
%deja definies ailleurs
\newcommand{\g}{\gamma}
\newcommand{\m}{\mu}
\newcommand{\G}{{\Gamma}}



%%%%%%%%%%%%%% pour les lettres calligraphi仔s

\newcommand{\Arond}{{\mathcal A}}
\newcommand{\Brond}{{\mathcal B}}
\newcommand{\Crond}{{\mathcal C}}
\newcommand{\Drond}{{\mathcal D}}
\newcommand{\Erond}{{\mathcal E}}
\newcommand{\Frond}{{\mathcal F}}
\newcommand{\Grond}{{\mathcal G}}
\newcommand{\Hrond}{{\mathcal H}}
\newcommand{\Irond}{{\mathcal I}}
\newcommand{\Jrond}{{\mathcal J}}
\newcommand{\Krond}{{\mathcal K}}
\newcommand{\Lrond}{{\mathcal L}}
\newcommand{\Mrond}{{\mathcal M}}
\newcommand{\Nrond}{{\mathcal N}}
\newcommand{\Orond}{{\mathcal O}}
\newcommand{\Prond}{{\mathcal P}}
\newcommand{\Qrond}{{\mathcal Q}}
\newcommand{\Rrond}{{\mathcal R}}
\newcommand{\Srond}{{\mathcal S}}
\newcommand{\Trond}{{\mathcal T}}
\newcommand{\Urond}{{\mathcal U}}
\newcommand{\Vrond}{{\mathcal V}}
\newcommand{\Wrond}{{\mathcal W}}
\newcommand{\Xrond}{{\mathcal X}}
\newcommand{\Yrond}{{\mathcal Y}}
\newcommand{\Zrond}{{\mathcal Z}}












 \def \qed {{\hspace*{\fill}$\halmos$\medskip}}
 \def \Z {{\mathbb Z}}
 \def \R {{\mathbb R}}
 \def \C {{\mathbb C}}
 \def \N {{\mathbb N}}
 \def \P {{\mathbb P}}
 \def \E {{\mathbb E}}
 \def \Q {{\mathbb Q}}
 \def \ra {\rightarrow}
 \def \da {\downarrow}
 \def \ua {\uparrow}

 \def \cO {{\mathcal O}}
 \def \cW {{\mathcal W}}
 \def \cD {{\mathcal D}}
 \def \cL {{\mathcal L}}
 \def \cI {{\mathcal I}}
 \def \cJ {{\mathcal J}}
 \def \cR {{\mathcal R}}
 \def \cA {{\mathcal A}}
 \def \cM {{\mathcal M}}
 \def \cN {{\mathcal N}}
 \def \cE {{\mathcal E}}
 \def \cP {{\mathcal P}}
 \def \cQ {{\mathcal Q}}
 \def \cB {{\mathcal B}}
 \def \cF {{\mathcal F}}
 \def \cG {{\mathcal G}}
 \def \cC {{\mathcal C}}
 \def \cV {{\mathcal V}}
 \def \cH {{\mathcal H}}
 \def \cT {{\mathcal T}}
 \def \cZ {{\mathcal Z}}
 \def \cS {{\mathcal S}}



 

\begin{document}

\textbf{Détermination d'une base explicite du noyau :}\\

1) Mettre la matrice sous forme échelonnée réduite.\\

\textit{On va considérer ici la matrice échelonnée réduite : }

$$\begin{pmatrix}
    1&2&0&-1&0\\
    0 &0&1&3&0\\
    0&0&0&0&1
\end{pmatrix}$$

2) Identifier les indices des colonnes sans pivot : chacune va donner un vecteur différent de la base. On note ces indices $j_1,...,j_r$.\\

\textit{Ici, ce sont les colonnes $j_1 = 2$ et $j_2 = 4$}\\

3) Prenons $k\in\{1,...,r\}$. On construit un vecteur de la base pour cet indice de la façon suivante :\\
- La $j_k$-ème coordonnée du vecteur est $1$. \\
- Toutes les coordonnées d'indices $j_p$ pour $p\neq k$ sont $0$.\\
- Toues les coordonnées $j $ pour $j> j_k$ sont $0$.\\
- Pour les coordonnées restantes (celles dont les indices correspondent aux colonnes avec un pivot, avec un indice plus petit que $j_k$), on place l'opposé des coefficients restants dans la colonne $j_k$. \\ 



\textit{Dans notre exemple, on commence avec $k = 1$ et $j_k = 2$. On construit un vecteur $u = (u_1,u_2,u_3,u_4,u_5)$. D'abord, puisque l'indice de la colonne est $j_k =2$, le 2ème coefficient vaut $1$ : $u = (u_1,1,u_3,u_4,u_5)$. Puis, on met des $0$ pour les indices correspondant aux autres colonnes sans pivot : $u = (u_1,1,u_3,0,u_5)$. Ensuite, les coefficients $u_3$ et $u_5$ sont nuls car $k = 2$, et $3> 2$ et $5 > 2$. Enfin, on complète avec l'opposé du coefficient restant sur la colonne $2$ : $u = (-2,1,0,0,0)$. 
}\\

Prenons maintenant $k = 2$, et donc $j_k = 4$. On construit $v = (v_1,v_2,v_3,v_4,v_5)$. On pose directement $v_4 = 1$ car c'est l'indice de la colonne, et $v_2 = 0$ car c'est une autre colonne sans pivot. $v_5 = 0$ car $5> 4$. Enfin, on place l'opposé des coefficients restants sur la colonne $4$ pour $v_1$ et $v_3$ qui correspondent au colonnes avec pivot restantes : $v_1 = 1$, et $v_3 = -3$. Cela donne : $v = (1,0,-3,1,0)$.\\

\textit{On vérifie que $u$ et $v$ sont bien dans le noyau de la matrice. \\}






Enfin, remarquons qu'on a déterminé ici une base du noyau de la forme échelonnée réduite de la matrice, mais le noyau d'une matrice et celui de sa forme échelonnée réduite coïncident, donc la base est commune aux deux !

\end{document}